% Based on: two accompanying sampling-based dynamics planning papers
% (whole-body non-prehensile manipulation under failures; liquid container
% manipulation with obstacle avoidance).

\chapter{Extending Sampling-Based Dynamics Planning to Constrained Settings}
\label{ch:constrained}

% TODO: Describe how the Poke-RRT line of thinking generalizes to two
% richer settings: constrained embodiments and constrained payloads.

\section{Whole-Body Non-Prehensile Manipulation Under Actuation Failures}
\label{sec:constrained:failures}

\subsection{Long Deployments and Locked-Joint Failures}
\label{sec:constrained:deployments}

% TODO: Motivate long deployments in harsh environments and the prevalence
% of locked-joint failures.

\subsection{Effects on Workspace, Manipulability, and Feasibility}
\label{sec:constrained:effects}

% TODO: Quantify how a locked joint reshapes the reachable workspace,
% manipulability, and feasible interactions.

\subsection{Whole-Body Non-Prehensile Manipulation as Recovery}
\label{sec:constrained:wholebody}

% TODO: Present whole-body non-prehensile manipulation as a recovery
% strategy for the failed embodiment.

\section{Liquid Container Manipulation with Obstacle Avoidance}
\label{sec:constrained:liquid}

\subsection{Task-Biased Geometric Planning}
\label{sec:constrained:biased}

% TODO: Geometric planner biased toward task-relevant configurations of
% the container.

\subsection{Analytical Static Models for Fill Level}
\label{sec:constrained:analytical}

% TODO: Tie end-effector orientation to fill level via analytical static
% models.

\subsection{Learned Dynamics Model for Trajectory Verification}
\label{sec:constrained:learned}

% TODO: Train a learned dynamics model on real-world data; use it to
% verify candidate trajectories.

\subsection{Where Analytical and Learned Components Meet}
\label{sec:constrained:hybrid}

% TODO: Discuss the boundary between analytical and learned model
% components in the planner.

\section{Lessons Learned}
\label{sec:constrained:lessons}

% TODO: Synthesize the strengths and limits of sampling combined with
% heuristics and analytical or learned models.

% TODO: Foreshadow the multi-query and learning-based directions taken in
% Chapters~\ref{ch:multiquery} through~\ref{ch:seeding}.
